\chapter{信息几何文献概述}\label{chap:introduction}

信息几何（Information Geometry）是利用微分几何方法研究概率分布空间的几何结构的学科。它为统计学、机器学习、热力学和量子力学等领域提供了统一的数学框架。

\section{什么是信息几何}

信息几何的核心思想是将概率分布族视为一个黎曼流形，利用几何方法来研究统计推断问题。这种观点使得我们能够：

\begin{itemize}
    \item 理解最大似然估计的几何意义
    \item 研究统计不变量和充分统计量
    \item 建立不同统计方法之间的联系
    \item 将信息论、统计学和微分几何有机结合
\end{itemize}

\section{核心概念}

\subsection{Fisher信息度量}

给定参数化概率分布族 $\{p(x;\theta)\}_{\theta \in \Theta}$，Fisher信息矩阵定义为

\begin{equation}
    g_{ij}(\theta) = \mathbb{E}\left[\frac{\partial \log p(X;\theta)}{\partial \theta_i} \cdot \frac{\partial \log p(X;\theta)}{\partial \theta_j}\right]
\end{equation}

Fisher信息度量是概率分布空间上的自然黎曼度量，具有统计学意义下的唯一性（Cramér-Rao定理）。

\subsection{对偶联络}

信息几何中的一个核心概念是对偶联络（dual connections）。对于Fisher度量 $g$，存在唯一的一对对偶联络 $(\nabla, \nabla^*)$ 满足

\begin{equation}
    Xg(Y,Z) = g(\nabla_X Y, Z) + g(Y, \nabla_X^* Z)
\end{equation}

这对偶联络对应于指数联络和混合联络，是信息几何许多理论的基础。

\section{文献目录}\label{sec:literature}

本学习资料涵盖以下文献，按主题分类：

\subsection{核心教材}

\begin{enumerate}
    \item \textbf{\textit{Information Geometry and Its Applications}} — Shun-ichi Amari

    这是信息几何领域的奠基之作，由该领域创始人编写。内容涵盖：
    \begin{itemize}
        \item 概率分布空间的几何结构
        \item Fisher信息与Cramér-Rao下界
        \item $\alpha$-联络与对偶平坦结构
        \item 指数族与混合族
        \item 统计学习中的几何方法
        \item 神经网络的自然梯度下降
        \item 信息几何在机器学习中的应用
    \end{itemize}

    \item \textbf{\textit{Information Geometry}} — Nihat Ay, Jürgen Jost, Hông Vân Lê, Lorenz Schwachhofer

    侧重于信息几何的数学基础，涵盖：
    \begin{itemize}
        \item 非参数统计流形
        \item 函数空间几何
        \item 进化博弈论中的应用
        \item 信息几何的泛函分析基础
    \end{itemize}

    \item \textbf{\textit{Elements of Information Theory}} (GTM 134) — Cover \& Thomas

    信息论经典教材，为理解信息几何中的KL散度等概念提供坚实基础：
    \begin{itemize}
        \item 熵与互信息
        \item KL散度与信息论不等式
        \item 信源编码与信道编码
        \item 信息论与统计推断的关系
    \end{itemize}

    \item \textbf{\textit{信息论基础}} — Cover \& Thomas（中文版）

    中文翻译版，便于参考。
\end{enumerate}

\subsection{非参数信息几何}

\begin{enumerate}
    \item \textbf{\textit{Nonparametric Information Geometry}} — Giovanni Pistone (arXiv:1306.0480)

    介绍了非参数统计流形的微分几何结构，是该领域的开创性工作。

    \item \textbf{\textit{A Lecture About the Use of Orlicz Spaces in Information Geometry}} — Giovanni Pistone (arXiv:2012.03376)

    2020年Les Houches暑期学校的讲义，介绍Orlicz空间在信息几何中的应用。

    \item \textbf{\textit{Connections on Non-Parametric Statistical Manifolds by Orlicz Space Geometry}} — Paolo Gibilisco \& Giovanni Pistone

    早期工作，利用Orlicz空间几何研究非参数统计流形上的联络。
\end{enumerate}

\subsection{Orlicz空间理论}

Orlicz空间是信息几何中处理非参数情况的重要工具：

\begin{enumerate}
    \item \textbf{\textit{Quantum Orlicz Spaces in Information Geometry}} — R.F. Streater (arXiv:math-ph/0407046)

    量子版本的Orlicz空间，用于量子信息几何。

    \item \textbf{\textit{On Musielak–Orlicz Function Spaces and Applications to Information Geometry}} — Rui Facundo Vigelis (PhD Thesis)

    博士学位论文，系统研究Musielak-Orlicz函数空间及其在信息几何中的应用。

    \item \textbf{\textit{New Results on Mixture and Exponential Models by Orlicz Spaces}} — Santacroce, Siri, Trivellato (arXiv:1603.05465)

    利用Orlicz空间方法研究混合模型和指数模型。

    \item \textbf{\textit{On Kōthe Duals of Orlicz-Lorentz Spaces}} — Kamińska \& Tag (arXiv:2404.07368)

    最新研究Orlicz-Lorentz空间的对偶性质。

    \item \textbf{\textit{When and Where the Orlicz and Luxemburg (quasi-) Norms are Equivalent?}} — del Campo et al.

    讨论Orlicz范数与Luxemburg范数的等价性条件。

    \item \textbf{\textit{Applications of Orlicz Spaces}} — M.M. Rao \& Z.D. Ren

    Orlicz空间的综合参考文献。

    \item \textbf{\textit{Topological Duals of Locally Convex Function Spaces}} — Pennanen \& Perkkiö (arXiv:2008.08934)

    研究局部凸函数空间的拓扑对偶，为信息几何提供泛函分析基础。
\end{enumerate}

\subsection{统计流形与几何结构}

\begin{enumerate}
    \item \textbf{\textit{Statistical Manifold as an Affine Space: A Functional Equation Approach}} — Zhang \& Hästö

    将统计流形视为仿射空间，利用函数方程方法研究其几何性质。

    \item \textbf{\textit{A Brief Introduction to Finsler Geometry}} — Matias Dahl

    Finsler几何简介，为理解信息几何中的更一般度量提供背景。
\end{enumerate}

\subsection{应用方向}

\begin{enumerate}
    \item \textbf{\textit{Information Geometry and Evolutionary Game Theory}} — Marc Harper (arXiv:0911.1383)

    进化博弈论中的Shahshahani几何与信息几何的关系。

    \item \textbf{\textit{Measuring Thermodynamic Length}} — Gavin E. Crooks

    热力学长度的信息几何视角，连接信息论与热力学。

    \item \textbf{\textit{Possible Generalization of Boltzmann-Gibbs Statistics}} — Constantino Tsallis

    Tsallis统计（广义Boltzmann-Gibbs统计）的原始论文，是非广延统计力学的基础。

    \item \textbf{\textit{Logit Standardization in Knowledge Distillation}} — Sun et al. (arXiv:2403.01427)

    信息几何在知识蒸馏中的应用。

    \item \textbf{\textit{Maximum Likelihood Estimation for the $\lambda$-exponential Family}} — Nielsen \& Barbaresco

    $\lambda$-指数族的最大似然估计，几何结构在统计推断中的应用。
\end{enumerate}

\subsection{高级主题}

\begin{enumerate}
    \item \textbf{\textit{Towards a Category-theoretic Foundation of Classical and Quantum Information Geometry}} — Ciaglia et al. (arXiv:2509.10262)

    信息几何的范畴论基础，最新的基础性工作。

    \item \textbf{\textit{On Reproducing Kernel Banach Spaces: Generic Definitions and Unified Framework of Constructions}} — Lin, Zhang, Zhang

    再生核Banach空间的统一构造框架，在信息几何中的应用潜力。
\end{enumerate}

\section{预备知识}

学习信息几何需要以下基础知识：

\begin{enumerate}
    \item \textbf{微分几何}：流形、切空间、协变导数、联络、黎曼几何
    \item \textbf{概率论}：概率分布、期望、方差、条件概率、极限定理
    \item \textbf{信息论}：熵、KL散度、互信息
    \item \textbf{统计学}：最大似然估计、充分统计量、Cramér-Rao下界
    \item \textbf{泛函分析}（高级主题）：Banach空间、Orlicz空间、函数空间几何
\end{itemize}

\section{学习路径建议}

建议按以下顺序学习：

\begin{enumerate}
    \item \textbf{入门}：Cover \& Thomas（信息论基础）$\rightarrow$ Amari（核心教材）
    \item \textbf{进阶}：Ay et al.（数学基础）$\rightarrow$ Pistone（非参数IG）
    \item \textbf{专题}：根据兴趣选择应用方向（进化博弈、热力学、量子IG等）
\end{enumerate}

\begin{Remark}
    本章为信息几何的导论章节，旨在建立整体认知框架。在后续章节中，我们将深入探讨各个核心概念。
\end{Remark}
